\lysh{13298}{4}

\begin{alist}
\item Αφού τα σημεία $A$ και $B$ είναι συμμετρικά ως προς τη διχοτόμο του $1$ου και του $3$ου τεταρτημορίου θα ισχύει $x_B = y_A$ και $y_B = x_A$.

\item Η κλίση της ευθείας $\varepsilon$ που διέρχεται από τα σημεία $A, B$ δίνεται από τη σχέση:
$\alpha = \dfrac{y_B - y_A}{x_B - x_A}$,
η οποία για τα $x_B = y_A$ και $y_B = x_A$, όπως προκύπτει από το ερώτημα $\alpha)$ γίνεται:
\[\alpha = \dfrac{x_A - y_A}{y_A - x_A} = -1.\]

\item 
\begin{rlist}
\item Από τη σχέση $y_A = x_B$ για $y_A = \kappa^2 - 3\kappa + 1$ και $x_B = \kappa - 2$, έχουμε:
\[\kappa^2 - 3\kappa + 1 = \kappa - 2 \Leftrightarrow \kappa^2 - 4\kappa + 3 = 0.\]
Η διακρίνουσα του τριωνύμου είναι $\Delta = (-4)^2 - 4 \cdot 3 = 4 > 0$ οπότε έχει δύο άνισες ρίζες:
$\kappa_1 = \dfrac{4 - 2}{2} = 1$ και $\kappa_2 = \dfrac{4 + 2}{2} = 3$.
Επειδή τα σημεία $A$ και $B$ ανήκουν στο 1ο τεταρτημόριο πρέπει $y_A > 0$ και $x_B > 0$.
Για $\kappa_1 = 1$ έχουμε $y_A = x_B = -1 < 0$, απορρίπτεται.
Για $\kappa_2 = 3$ έχουμε $y_A = x_B = 1 > 0$, δεκτή.
Άρα $\kappa = 3$.
Οπότε, τα σημεία $A$ και $B$ έχουν συντεταγμένες $(4, 1)$ και $(1, 4)$ αντίστοιχα.
\item Η ευθεία $\varepsilon$ έχει εξίσωση της μορφής $y = \alpha x + \beta$ και κλίση $\alpha = -1$, άρα
$y = -x + \beta$.
Επειδή διέρχεται από το σημείο $A$ που για $\kappa = 3$ έχει συντεταγμένες $(4, 1)$, θα ισχύει:
\[1 = -4 + \beta \Leftrightarrow \beta = 5.\]
Άρα η εξίσωση της $\varepsilon$ είναι η $y = -x + 5$.
\end{rlist}
\end{alist}
